\chapter{相关工作与技术}

\section{强化学习基础}
强化学习是一种通过智能体与环境的交互来学习最优行为策略的机器学习方法\cite{citylearn2020}。智能体、动作、环境、状态和奖励在强化学习中相互关联、相互作用，共同构成了一个完整的学习系统。建筑群能源控制可以抽象为马尔可夫决策过程$(S,A,P,R,\gamma)$：状态$s_t\in S$包含气象、电价、负荷和设备信息，动作$a_t\in A$为设备设定值，转移$P(s_{t+1}|s_t,a_t)$由建筑热动态与设备物理决定，奖励$r_t$反映成本与舒适性。强化学习目标是寻找策略$\pi(a|s)$最大化期望折扣回报，如式\ref{eq:rl-obj}所示。
\begin{equation}
\label{eq:rl-obj}
J(\pi)=\mathbb{E}_{\pi}\left[\sum_{t=0}^{T}\gamma^t r(s_t,a_t)\right].
\end{equation}

（1）智能体：在建筑群场景中，每栋建筑或每类设备被建模为独立智能体。智能体观察周遭环境的状态变化，选择并执行相应动作，并根据环境反馈的奖励或惩罚调整决策策略，追求长期累积奖励最大化\cite{marl2021}。

（2）动作：在本文中动作为CityLearn设备设定值向量，对应HVAC温度设定、DHW功率与电池充放电设定。动作的选择是智能体策略的具体体现，不同策略导致不同的负荷曲线与社区峰值。

（3）环境：环境是智能体所处的建筑群物理世界，接受动作输入并产生新的状态与奖励反馈。CityLearn通过Gymnasium接口提供标准化的气象、电价、碳强度与负荷时序，成为多智能体需求响应的公共仿真场所\cite{citylearn2020}。

（4）状态：状态是环境在某一时刻的完整描述。在本文中建筑观测包括室外温度、太阳辐照、电价、碳强度、建筑电/热负荷、室内温度与电池SOC；社区汇总特征包括总负荷、平均SOC与峰值指示。状态可以是离散的，也可以是连续的传感器与时序数据。

\subsection{基于值函数的强化学习算法}
强化学习算法在策略更新和学习方法上可明确划分为两大类：基于值函数的强化学习算法和基于直接策略搜索的强化学习算法。基于值函数的方法核心是学习状态值函数或动作值函数，不断迭代更新直到收敛，再通过贪婪策略选择使值函数最大的动作。在建筑群控制中，DQN用神经网络近似$Q$函数并用经验回放稳定训练\cite{mnih2015}；DDPG针对连续动作引入确定性策略与评论家网络\cite{ddpg2016}。但值函数方法在多建筑连续设定值空间中容易面临维度灾难，需要结合策略搜索或残差结构降低探索复杂度。

\subsection{基于直接策略搜索的强化学习算法}
基于直接策略搜索的方法直接学习从状态到动作分布的映射，通过策略梯度或信赖域方法更新参数。PPO通过裁剪目标限制策略更新步长，提高了训练稳定性\cite{ppo2017}，是本文独立多智能体基线与残差策略训练的参考方法。直接策略搜索天然适合电池充放电、HVAC温度设定等连续动作，但纯策略探索可能产生违反设备约束或舒适性要求的动作，因此本文将其与CHESCA基线结合，仅学习小范围残差修正。

\subsection{多智能体扩展与经验池}
当每栋建筑视为智能体$i\in\{1,\dots,N\}$时，问题变为马尔可夫博弈，联合观测为$\mathbf{o}=(o_1,\dots,o_N)$，联合动作为$\mathbf{a}=(a_1,\dots,a_N)$。主要难点包括环境非平稳性、信用分配与通信约束\cite{marl2021}。集中训练分散执行（CTDE）是主流范式\cite{lowe2017}：训练阶段评论家可访问全局信息，执行阶段策略仅依赖本地观测。MADDPG为每个智能体维护Actor与集中式Critic\cite{lowe2017}；Wilk等进一步提出仅共享离散化``策略意图''而非全部观测的LSD-MADDPG，在保护隐私的同时保持协同\cite{wilk2024}。经验池通过存储历史交互样本并随机采样，打破时序相关性，是稳定多智能体训练的重要机制；本文训练入口据此保存观测、基线动作、残差与KPI增量，供离线复训与消融分析。

\section{可视化技术}
可视化技术作为一种直观高效的信息传递手段，在建筑群能源实验的管理和优化中发挥着不可替代的作用。通过将复杂的负荷时序、KPI权衡与决策过程以图形形式呈现，可视化帮助研究人员快速把握整体性能与局部细节，发现潜在问题与优化方向。
\subsection{Vue}
前端采用Vue与Element UI构建单页应用。Vue的组件化机制将系统首页、模型配置、能源分析与决策轨迹封装为独立组件，通过路由与状态管理组织页面跳转与任务轮询；Element UI提供表单、表格、对话框等基础控件，保证配置页面与任务列表的交互一致性。平台首页聚合社区KPI卡片，模型优选页支持多算法多指标对比，更多信息与决策对话框支持逐步查看基线动作、残差值与掩码记录。

\subsection{Echarts}
图表层采用Apache ECharts库构建，确保数据可视化既美观又高效。页面中心根据展示内容划分为多个独立而相互关联的区域：折线图直观呈现负荷、电价、SOC与室温随时间步变化的趋势；柱状图对比不同算法的KPI；雷达图揭示成本、碳排、峰值、爬坡与舒适性之间的权衡；决策轨迹图同步展示基线动作、残差与投影结果。能源语义色在全平台保持一致：电网使用蓝色、光伏使用橙色、电池使用青绿色。

\section{多智能体强化学习与残差控制}
\subsection{残差强化学习}
残差强化学习将最终动作分解为基线与修正\cite{residual2019}，如式\ref{eq:residual}所示：
\begin{equation}
\label{eq:residual}
a_t = a^{base}_t + \alpha \cdot \Delta a^{RL}_t,\quad \Delta a^{RL}_t\sim\pi_\theta(\cdot|o_t),
\end{equation}
其中$\alpha$为残差幅度系数。基线$a^{base}_t$由CHESCA、RBC或规划器给出，保证基本可行性；策略仅需学习小范围修正，从而降低探索空间与约束违反概率。在机器人控制中该方法显著提高了样本效率\cite{residual2019}。在能源领域，残差结构天然对应``规则保安全、学习求增益''：若策略输出为零则退化为基线；安全投影层再将动作裁剪到物理边界内。本文将残差系数$\alpha$、动作掩码与安全后处理作为可配置项，与CityLearn动作接口解耦。

\subsection{CHESCA社区分层协调思想}
2023年CityLearn Challenge的优胜方案CHESCA采用社区分层结构\cite{chesca2024}。其要点可概括为：日前层利用XGBoost预测未来负荷并进行电价感知规划；社区层计算总负荷修正信号并分配各建筑参考；本地层跟踪参考并通过路径搜索优化电池充放电序列，成本函数综合电价、碳与爬坡惩罚。该方法仅用有限历史数据即取得稳定性能，体现了``规则先验+数据驱动''的优势。CHESCA的局限是参数静态，难以适应负荷漂移；本文在其上叠加ResMARL，正是为了保留其稳定性并引入自适应修正能力。

\subsection{CityLearn仿真环境与评价指标}
CityLearn是基于Gymnasium接口的开源建筑群仿真平台，专为分布式能源资源调度设计\cite{citylearn2020}。其核心贡献是统一任务框架：提供历史气象、电价、碳强度与负荷时序，定义HVAC、DHW与电池动作接口，并给出多目标KPI\cite{citylearn2024}。2020年后每年举办的CityLearn Challenge进一步推动算法在同一数据与指标下比较\cite{citylearnchallenge2023}。CityLearn v2扩展了韧性、住户中心与碳感知能力，支持电网交互式社区管理\cite{citylearn2024}。本文采用的KPI包括用电成本、碳排放、不舒适度、峰值、爬坡、可再生消纳与电网净负荷等。评价时通常以无控制（NOCONTROL）为分母做归一化，以便跨场景比较。BOPTEST等基于Modelica的框架虽物理精度更高，但部署复杂，更适合单体建筑控制器测试\cite{boptest2021}；CityLearn更适合社区级策略快速迭代，因而被本文选为基准。

\section{平台技术与数据管理}
后端采用Spring Boot提供RESTful接口，MySQL持久化用户、社区、建筑、算法任务与KPI，Redis承担任务队列与状态缓存；Python算法组件独立进程运行，通过文件与任务目录与Java交换配置和结果；前后端分离架构使长时间仿真任务可异步执行，前端轮询任务状态并展示中间日志。数据层面，数据库以建筑时序表、环境时序表、算法脚本表、任务表、KPI结果表、算法配置表与数据集目录表组织CityLearn输入、任务状态与评价结果，详细表设计见第三章概念设计与第四章4.3节。

\section{本章小结}
本章介绍了为了深入研究建筑群节能控制与多智能体残差策略，所应具备的理论知识。通过对研究对象——强化学习基础、值函数与策略搜索方法、多智能体扩展的了解，对可视化技术Vue与ECharts的进一步探索，以及对残差控制、CHESCA分层思想、CityLearn环境与平台数据管理的全面总结，为本文后续系统设计、算法实现与可视化管理奠定了理论基础和技术导向。
